\documentclass[11pt,a4paper]{article}

\usepackage[spanish]{babel}


\usepackage{ifthen}
\usepackage{amsmath}
\usepackage[T1]{fontenc}
%\usepackage{lucidabr}
\usepackage{alltt}
\usepackage{pifont}
\usepackage[dvipsnames]{xcolor}
\usepackage{graphicx}
\usepackage{pst-all}

\usepackage[verbose]{geometry}
\geometry{a4paper}
\geometry{margin=1.7cm,nohead}


\def\pathFM{/home/alberto/famosoLibro/}


%\input{\pathFM algoritmos}  % comandos del pseudo-lenguaje
%\input{\pathFM especificaciones}  % 

% pseudo-codigo para escribir algoritmos

% tipo se letra para las palabras reservadas
\newcommand{\reser}{\textbf}
%  letra para los identificadores de procedimiento y funcion
\newcommand{\T}[1]{\textsf{#1}}
%  letra para los tipos 
\newcommand{\ti}[1]{\textit{\textsf{#1}}}
%  letra para los identificadores de variables  
\newcommand{\V}[1]{\mbox{\textit{#1}}}
%nombres de operaciones
\newcommand{\op}[1]{\textsf{#1}}

\newcommand{\nat}{\mathrm{I}\!\mathrm{N}}
\newcommand{\real}{\mathrm{I}\!\mathrm{R}}

% entorno de algoritmos
% \newcommand{\alg}{\medskip\begin{small}\noindent$}
% \newcommand{\falg}{$\end{small}\medskip}
\newcommand{\alg}{\medskip\begin{small}\noindent$}
\newcommand{\falg}{$\end{small}\medskip}

% comentarios
\newcommand{\COMMENT}[1]{\{\;\mbox{#1}\;\}}

% palabras reservadas
\newcommand{\fun}{\reser{fun} \;\,}
\newcommand{\ffun}{\reser{ffun}}
\newcommand{\proc}{\reser{proc} \;\,}
\newcommand{\fproc}{\reser{fproc}}
\newcommand{\e}{\reser{e}\;}
\newcommand{\s}{\reser{s}\;}
\newcommand{\var}{\reser{var} \;\,}
\newcommand{\dev}{\; \, \reser{dev} \;\,}
\newcommand{\para}{\reser{para} \;\,}
\newcommand{\hasta}{\;\, \reser{hasta} \;\,}
\newcommand{\paso}{\;\, \reser{paso} \;\,}
\newcommand{\fpara}{\reser{fpara}}
\newcommand{\mientras}{\reser{mientras} \;\,}
\newcommand{\hacer}{\;\, \reser{hacer} \;\,}
\newcommand{\fmientras}{\reser{fmientras}}
\newcommand{\repetir}{\reser{repetir} \;\,}
\newcommand{\frepetir}{\reser{frepetir}}
\newcommand{\hastar}{\reser{hasta} \;\,}
\newcommand{\sip}{\reser{si} \;\, }
\newcommand{\entonces}{\;\, \reser{entonces} \;\,}
\newcommand{\sino}{\reser{si no} \;\, }
\newcommand{\fsi}{\reser{fsi}}
\newcommand{\con}{\reser{con} \;\, }
\newcommand{\fcon}{\reser{fcon} \;\, }
\newcommand{\casos}{\reser{casos}}
%\newcommand{\oo}{\pspicture(1em,0.8em)%
%\psframe[linewidth=.5pt,unit=1em](0.4,0.8)\endpspicture}
%\newcommand{\ooV}{\pspicture(1em,1em)\endpspicture}
%\newcommand{\ooV}{\phantom{\makebox[0pt][r]{[\!]\ }}}
%\newcommand{\oo}{\makebox[0pt][r]{[\!]\ }}
\newcommand{\oo}{\scalebox{0.7}[1.4]{\makebox[0pt]{$\sqcup$}\makebox[0pt]{$\sqcap$}\ \ \ }}
\newcommand{\ooV}{\phantom{\oo}}
\newcommand{\fl}{\; \rightarrow \; \;}
\newcommand{\fcasos}{\reser{fcasos}}
\newcommand{\nada}{\reser{nada}}
\newcommand{\de}{\;\reser{de}\;}
\newcommand{\reg}{\reser{reg}}
\newcommand{\freg}{\reser{freg} }
\newcommand{\punteroa}{\reser{puntero a}}

% div y mod en definiciones matematicas
\renewcommand{\div}{\; \mathrm{div} \;} 
%\newcommand{\mod}{\; \mathrm{mod} \;}
% div y mod en algoritmos
\newcommand{\Tdiv}{\; \T{div} \;}
\newcommand{\Tmod}{\; \T{mod} \;}

% min y max en algoritmos
\newcommand{\Tmin}{\T{m\'\i n}}
\newcommand{\Tmax}{\T{m\'ax}}

\newcommand{\nil}{\T{nil}}
\newcommand{\conc}{\mathbin{+\!\!+}}

\newcommand{\tipo}{\reser{tipos}}
\newcommand{\ftipo}{\reser{ftipos}}

% booleanos
\newcommand{\bool}{\ti{bool}}
\newcommand{\cierto}{\textsf{cierto}}
\newcommand{\falso}{\textsf{falso}}

% otros simbolos
\newcommand{\asig}{\, : = \,}  % asignacion
\newcommand{\asigM}[1]{\langle \, #1 \, \rangle} % para la asignacion multiple
\newcommand{\result}[1]{\langle \, #1 \, \rangle} % para los resultados
\newcommand{\tupla}[1]{\langle\,#1\,\rangle} % para las tuplas
\newcommand{\seq}{\; ; \;\;} % punto y coma (instruccion secuencial)
\newcommand{\yl}{\;\wedge\;} % y logica
\newcommand{\ylc}{\;\wedge_c\;} % y logica con cortocircuito
\newcommand{\ol}{\;\vee\;}   % o logica
\newcommand{\olc}{\;\vee_c\;}   % o logica con cortocircuito
\newcommand{\nol}{\neg}  % no logica
\newcommand{\punt}{\!\uparrow\!} % puntero

% espacios para indentar los algoritmos
\newcommand{\uno}{\hspace*{1.5em}}
\newcommand{\dos}{\hspace*{3em}}
\newcommand{\tres}{\hspace*{4.5em}}
\newcommand{\cuatro}{\hspace*{6em}}
\newcommand{\cinco}{\hspace*{7.5em}}
\newcommand{\seis}{\hspace*{9em}}
\newcommand{\siete}{\hspace*{10.5em}}
\newcommand{\ocho}{\hspace*{12em}}
\newcommand{\nueve}{\hspace*{13.5em}}


\newcommand{\reservar}{\T{reservar}}
\newcommand{\liberar}{\T{liberar}}


% booleanos para controlar las listas de ejercicios

% para decir si se muestran los enunciados o no
\newboolean{conEnunciado}
\setboolean{conEnunciado}{true} % por defecto

% para decir si se muestran las soluciones o no
\newboolean{solucion}
\setboolean{solucion}{false} % por defecto

% contadores
\newcounter{cont}
\newcounter{exercises}
\def\theexercises{\arabic{exercises}}

\newenvironment{listaejercicios}
{ \ifthenelse{\boolean{solucion}}
  {\begin{list}{\refstepcounter{exercises}%
  \addtocounter{exercises}{-1}%
                \textbf{Ejercicio \theexercises\enspace}}
                {\usecounter{exercises}
                 \setlength{\labelwidth}{0cm}
                \setlength{\labelsep}{0cm}
                \setlength{\itemindent}{0cm}
                \setlength{\listparindent}{\parindent}
                \setlength{\leftmargin}{0cm}
                 }
  }
  {\begin{list} {\refstepcounter{exercises}%
  \addtocounter{exercises}{-1}%
                \textbf{Ejercicio \theexercises\enspace}}%
  {\usecounter{exercises}
  }
  }         
}
{\end{list}}

\newcommand{\cuestion}{\addtocounter{exercises}{1}\textbf{Cuesti\'on \arabic{exercises}.}}

\newcommand{\hecho}{\makebox[0pt]{\hskip-0.8cm\psscalebox{-1.0 1}{\Large\ding{45}}}}
%\newcommand{\pepe}{\makebox[0pt]{\hskip-1.3cm\raisebox{-.15cm}{\scalebox{.05}{\includegraphics{gorro.eps}}}}}
\newcommand{\enLibro}{\item[\addtocounter{exercises}{1}%
\hecho\textbf{Ejercicio \arabic{exercises}\enspace}]}

\newcommand{\Label}[1]{%
\ifthenelse{\boolean{conEnunciado}}{\label{#1}}{}}

% \newcommand{\enunciado}[1]{
%   \ifthenelse{\boolean{conEnunciado}}
%   {  \ifthenelse{\boolean{solucion}}
%         % con soluciones
%         {~\textsf{#1}}
%         % sin soluciones
%         {#1}
%   }
%   {}}

\newcommand{\enunciado}[1]{#1}

\newcommand{\solucion}[1]{%
\ifthenelse{\boolean{solucion}}{\hazsolucion{#1}}{}}

\newcommand{\hazsolucion}[1]{ 
 \ifthenelse{\boolean{conEnunciado}}{}{\mbox{}\\[-0.3cm]}
%\begin{list}{}{\setlength{\leftmargin}{0cm}}
%\item[]
 \ifthenelse{\boolean{conEnunciado}}{\noindent\textbf{SOLUCI\'ON:}\\}{}

#1
%\end{list}
%\vspace{1.5cm}
% la raya que pinta debajo
% \begin{center}
% \psset{unit=1cm}
% \begin{pspicture}(0,0.5)(10,1.5)
% \psframe*(0,1)(10,1.1)
% \end{pspicture}
% \end{center}
\begin{center}
$\Box$%\flushright \Large{\Peace}\\
\end{center}
}
%\renewcommand{\hazsolucion}[1]{#1}

\newcounter{contAlph}
\newenvironment{listaalph}{\begin{list}{(\alph{contAlph})}%
                          {\usecounter{contAlph}\leftmargin 2em%
                            \setlength{\labelsep}{1ex} } }
                          {\end{list}}

\renewenvironment{enumerate}{\begin{list}{\arabic{cont}.}
                            {\usecounter{cont}\leftmargin 2em%
                            \setlength{\labelsep}{1ex} }}
                          {\end{list}}

\renewenvironment{itemize}{\begin{list}{$\bullet$}
                           {\usecounter{cont}\leftmargin 2em%
                           \setlength{\labelsep}{1ex}}}
                          {\end{list}}
 
\newenvironment{bibliografia}{\begin{small}%
\vspace*{1cm}{\noindent\large\bf Bibliograf\'{\i}a}
\begin{list}{[\arabic{cont}]}
                            {\usecounter{cont}\leftmargin 2em%
                            \setlength{\labelsep}{1ex} }}
                          {\end{list}\end{small}}

\newcommand{\Apartado}[1]{%
\subsubsection*{Apartado (#1)}%
}
\newcommand{\miCaption}[1]{%
\caption{#1}
}
\newenvironment{Figure}[1][tbp]{%
\begin{figure}[#1]
}{%
\end{figure}
}

\newcommand{\falta}[1]{\vspace*{1cm} \begin{center}\psshadowbox{Falta #1}%
                       \end{center}\vspace*{1cm}}
\newcommand{\NOTA}[1]{\psframebox[framesep=6pt]{\parbox{12cm}{\textsl{#1}}}}


\newcommand{\ds}{\displaystyle}

% para los dibujos con arboles de ramificacion y poda
\newcommand{\stack}[2]{\begin{array}{c} #1 \\ #2 \end{array}}
\newcommand{\stackT}[3]{\begin{array}{c} #1 \\ #2 \\ #3 \end{array}}


\newcommand{\examen}[1]{}


\newcommand{\cabecera}{%
\begin{center}
{\Large Facultad de Inform\'atica}\\[2ex]
{\Large \bf Metodolog\'\i a y Tecnolog\'\i a de la Programaci\'on}\\[2ex]
{\large Superior -- PLAN 98 -- curso 00-01}
\end{center}
}%

\newcommand{\titulo}[1]{%
% \vspace{1em}
\medskip
\begin{center}
{\Large \bf #1}
\end{center}
% \vspace{1em}
%{\Large \bf #1}
}%




\usepackage{url}
\usepackage{hyperref}
\usepackage{array}
\usepackage{enumitem}

\newcommand{\cb}{\ding{113}}
\newcommand{\cbfive}{\cb & \cb & \cb & \cb & \cb}
\newcommand{\cbfour}{\cb & \cb & \cb & \cb}

\newcommand{\bloque}[1]{\bigskip\noindent\textbf{\large #1}\par\nobreak\smallskip}

\begin{document}

\begin{center}
{\Large\textbf{Auditor\'ia, Calidad y Fiabilidad Inform\'aticas}}
\smallskip

\large
Universidad Complutense de Madrid \qquad Curso 2026/27
\medskip

Encuesta final sobre el proyecto de Maude\hfill\mbox{}

\noindent\hrulefill
\end{center}

\noindent
\textbf{Antes de empezar.} Esta encuesta es \textbf{an\'onima} y
\textbf{voluntaria}: no puntúa, no se sabe qui\'en la ha contestado y no
influye en tu nota ni en la de nadie.
%
Se tarda unos 10 minutos.
%
Este a\~no la asignatura se ha evaluado de una forma distinta a la habitual, y
lo que respond\'ais aqu\'i es lo \'unico que permite saber si el cambio ha
merecido la pena o hay que deshacerlo. Contestad con calma y sin adornar
nada: una respuesta cr\'itica es m\'as \'util que una amable.

\smallskip
\noindent
\textbf{Uso de los datos.} Las respuestas podr\'an publicarse en un art\'iculo
sobre el rediseño de la asignatura, incluidas respuestas individuales y citas
literales de las preguntas abiertas. No se publicar\'a tu nombre ni ning\'un
dato que permita identificarte. Si prefieres que tus respuestas no se usen
con ese fin, marca la casilla del \'ultimo apartado.

\bloque{A. Contexto}

\begin{enumerate}[leftmargin=*]

\item ¿De cu\'antas personas era tu equipo?\\
\cb\ Dos \qquad \cb\ Tres

\item El sistema que modelasteis lo elegisteis:\\
\cb\ De la lista del enunciado \qquad
\cb\ Lo propusimos nosotros

\item ¿Cu\'anto conoc\'ias ese sistema \emph{antes} de empezar el proyecto?\\
\cb\ Nada \quad
\cb\ Hab\'ia o\'ido hablar de \'el \quad
\cb\ Lo hab\'ia usado o estudiado por encima \quad
\cb\ Lo conoc\'ia bien

\item \emph{(Opcional)} ¿Qu\'e sistema modelasteis?
\rule{7cm}{0.4pt}\\
{\small Puedes dejarlo en blanco: si tu equipo fue el \'unico que eligi\'o
ese sistema, responder aqu\'i permite identificarlo (aunque no a ti dentro
de \'el).}

\end{enumerate}

\bloque{B. Dedicaci\'on}

\begin{enumerate}[leftmargin=*,resume]

\item Aproximadamente, ¿cu\'antas horas \textbf{en total} dedicaste al
proyecto fuera de clase?\\
\cb\ Menos de 10 \quad
\cb\ 10--19 \quad
\cb\ 20--29 \quad
\cb\ 30--39 \quad
\cb\ 40 o m\'as

\item ¿C\'omo se reparti\'o aproximadamente ese tiempo? (indica un porcentaje
orientativo en cada fila; deber\'ia sumar 100)

\smallskip
\begin{tabular}{@{}p{9cm}l@{}}
Entender el sistema real (documentaci\'on, art\'iculos, c\'odigo) & \rule{1.5cm}{0.4pt}\,\% \\
Escribir y depurar la especificaci\'on en Maude & \rule{1.5cm}{0.4pt}\,\% \\
Formular y verificar las propiedades (\texttt{search}, LTL) & \rule{1.5cm}{0.4pt}\,\% \\
Mutaciones & \rule{1.5cm}{0.4pt}\,\% \\
Redactar la memoria y preparar la defensa & \rule{1.5cm}{0.4pt}\,\% \\
\end{tabular}

\item Comparado con lo que esperabas al leer el enunciado, la carga de
trabajo real fue:\\
\cb\ Mucho menor \quad
\cb\ Menor \quad
\cb\ La esperada \quad
\cb\ Mayor \quad
\cb\ Mucho mayor

\end{enumerate}

\bloque{C. Aprendizaje percibido}

\noindent
Indica tu grado de acuerdo con cada afirmaci\'on:
{\small 1 = totalmente en desacuerdo, 2 = en desacuerdo, 3 = ni de acuerdo ni
en desacuerdo, 4 = de acuerdo, 5 = totalmente de acuerdo.}

\smallskip
\begin{tabular}{@{}p{11.2cm}ccccc@{}}
& \small 1 & \small 2 & \small 3 & \small 4 & \small 5 \\
\hline
8. El proyecto me hizo conectar ideas de la asignatura con conocimientos de otras asignaturas. & \cbfive \\[2pt]
9. Tuve que analizar en profundidad una idea o un mecanismo, no solo aplicar una receta. & \cbfive \\[2pt]
10. He aprendido algo que cambi\'o mi forma de entender c\'omo funcionan los sistemas concurrentes reales. & \cbfive \\[2pt]
11. Ahora s\'e decidir mejor cu\'ando algo debe ser una ecuaci\'on y cu\'ando una regla. & \cbfive \\[2pt]
12. Ahora entiendo mejor qu\'e garantiza y qu\'e \emph{no} garantiza una verificaci\'on formal. & \cbfive \\[2pt]
13. Sabr\'ia aplicar lo aprendido a un sistema distinto del que model\'o mi equipo. & \cbfive \\[2pt]
14. Decidir c\'omo representar el sistema (qu\'e tipos y constructores usar, qu\'e es estado y qu\'e es mensaje) me cost\'o m\'as que escribir las ecuaciones y las reglas. & \cbfive \\[2pt]
\hline
\end{tabular}

\bloque{D. Implicaci\'on con el proyecto}

\noindent
Para cada frase, indica en qu\'e medida te describe:
{\small 1 = no me describe nada, 2 = me describe poco, 3 = me describe
bastante, 4 = me describe totalmente.}

\smallskip
\begin{tabular}{@{}p{12.1cm}cccc@{}}
& \small 1 & \small 2 & \small 3 & \small 4 \\
\hline
15. En las sesiones dedicadas al proyecto trabajaba concentrado. & \cbfour \\[2pt]
16. Dediqu\'e al proyecto m\'as esfuerzo del m\'inimo necesario para aprobarlo. & \cbfour \\[2pt]
17. Disfrut\'e trabajando en el proyecto. & \cbfour \\[2pt]
18. Cuando algo no sal\'ia, me apetec\'ia insistir hasta entender por qu\'e. & \cbfour \\[2pt]
19. El proyecto me pareci\'o aburrido. & \cbfour \\[2pt]
20. Me limit\'e a hacer lo justo para cumplir con el enunciado. & \cbfour \\[2pt]
21. He hablado del proyecto con gente de fuera de la asignatura. & \cbfour \\[2pt]
22. Me qued\'e con ganas de seguir mejorando el modelo despu\'es de la entrega. & \cbfour \\[2pt]
\hline
\end{tabular}

\bloque{E. Comparaci\'on con las pr\'acticas cortas del principio de curso}

\noindent
Antes del proyecto hicisteis tres pr\'acticas cortas de Maude en esta misma
asignatura, que no contaban para la nota. Comparado con aquellas
pr\'acticas, el proyecto me result\'o:

\noindent
{\small 1 = mucho menos, 2 = algo menos, 3 = igual, 4 = algo m\'as,
5 = mucho m\'as.}

\smallskip
\begin{tabular}{@{}p{11.2cm}ccccc@{}}
& \small 1 & \small 2 & \small 3 & \small 4 & \small 5 \\
\hline
23. \dots interesante & \cbfive \\[2pt]
24. \dots exigente & \cbfive \\[2pt]
25. \dots \'util para aprender Maude de verdad & \cbfive \\[2pt]
26. \dots relacionado con lo que espero hacer profesionalmente & \cbfive \\[2pt]
27. \dots exigente en cuanto a entender el sistema \emph{antes} de poder
escribir c\'odigo & \cbfive \\[2pt]
\hline
\end{tabular}

\smallskip
\begin{enumerate}[leftmargin=*,resume]
\setcounter{enumi}{27}
\item Si pudieras elegir el formato de evaluaci\'on de la asignatura, preferir\'ias:\\
\cb\ Solo pr\'acticas cortas como las del principio de curso \quad
\cb\ Solo un proyecto en equipo como este \quad\\
\cb\ Las dos cosas, como este a\~no \quad
\cb\ Me da igual
\end{enumerate}

\bloque{F. Asistentes de IA y el MCP para Maude}

\noindent
El uso de asistentes de IA en este proyecto estaba \textbf{permitido}, y el
MCP para Maude recomendado, as\'i que aqu\'i no se trata de detectar nada: interesa saber c\'omo os han
funcionado y, sobre todo, si el MCP para Maude aporta algo que un asistente
gen\'erico no d\'e. El bloque es an\'onimo, como el resto.

\begin{enumerate}[leftmargin=*,resume]

\item ¿Con qu\'e frecuencia usaste alg\'un asistente durante el proyecto?\\
\cb\ Nunca \quad
\cb\ Puntualmente \quad
\cb\ Con frecuencia \quad
\cb\ Casi continuamente

\item ¿Cu\'ales? \emph{(marca todo lo que corresponda)}\\
\cb\ ChatGPT \quad \cb\ Claude \quad \cb\ Copilot \quad \cb\ Gemini \quad
\cb\ Otro: \rule{4cm}{0.4pt}

\item ¿Para qu\'e los usaste? \emph{(marca todo lo que corresponda)}

\begin{itemize}[leftmargin=*,label=\cb,itemsep=1pt]
\item Entender el sistema real que \'ibamos a modelar
\item Consultar sintaxis de Maude
\item Explicar mensajes de error del int\'erprete
\item Escribir ecuaciones o reglas concretas
\item Escribir la especificaci\'on entera o casi entera
\item Formular las propiedades (\texttt{search}, LTL)
\item Interpretar la salida del model checker
\item Proponer las mutaciones
\item Redactar partes de la memoria
\item Otros: \rule{5cm}{0.4pt}
\end{itemize}

\item ¿Usaste el MCP para Maude (Harold)?\\
\cb\ No lo instal\'e \quad
\cb\ Lo instal\'e pero apenas lo us\'e \quad
\cb\ Lo us\'e a veces \quad
\cb\ Lo us\'e habitualmente

\smallskip
{\small Si no lo usaste, indica brevemente por qu\'e y pasa a la
pregunta~\ref{q:rehacer}: \rule{5.5cm}{0.4pt}}

\item Instalarlo y ponerlo en marcha fue sencillo:\\
\cb\ 1 (nada) \quad \cb\ 2 \quad \cb\ 3 \quad \cb\ 4 \quad \cb\ 5 (mucho)

\item Los diagn\'osticos que devuelve sobre la especificaci\'on me resultaron
\'utiles:\\
\cb\ 1 (nada) \quad \cb\ 2 \quad \cb\ 3 \quad \cb\ 4 \quad \cb\ 5 (mucho)

\item Con el MCP, el asistente generaba c\'odigo Maude m\'as correcto que sin
\'el:\\
\cb\ 1 (nada) \quad \cb\ 2 \quad \cb\ 3 \quad \cb\ 4 \quad \cb\ 5 (mucho)
\quad \cb\ No lo puedo comparar

\item ¿En qu\'e situaciones te result\'o m\'as \'util el MCP?\\[8pt]
\rule{\linewidth}{0.4pt}

\item ¿Qu\'e le falta al MCP o qu\'e le cambiar\'ias?\\[8pt]
\rule{\linewidth}{0.4pt}\\[10pt]
\rule{\linewidth}{0.4pt}

\item\label{q:rehacer} Del c\'odigo Maude que te gener\'o un asistente, ¿cu\'anto tuviste que
rehacer o corregir?\\
\cb\ No us\'e asistentes para el c\'odigo \quad \cb\ Nada o casi nada \quad
\cb\ Poco \quad \cb\ Bastante \quad \cb\ Casi todo

\item ¿Entend\'ias todo el c\'odigo que entregasteis, viniera de donde
viniera?\\
\cb\ S\'i, todo \quad \cb\ Casi todo \quad \cb\ Hab\'ia partes que no acababa de entender \quad
\cb\ Bastantes partes no las entend\'ia

\item ¿Probasteis a pedirle a un asistente el proyecto entero a partir del
enunciado, a ver qu\'e sal\'ia?\\
\cb\ No lo probamos \quad
\cb\ S\'i, y no dio nada aprovechable \quad
\cb\ S\'i, y dio algo que hubo que rehacer casi entero \quad
\cb\ S\'i, y funcion\'o razonablemente bien

\item Comparado con las pr\'acticas cortas del principio de curso, resolver
este proyecto con ayuda de una IA me pareci\'o:\\
\cb\ Mucho m\'as f\'acil \quad \cb\ M\'as f\'acil \quad \cb\ Igual \quad
\cb\ M\'as dif\'icil \quad \cb\ Mucho m\'as dif\'icil

\item ¿Qu\'e es lo que peor hacen los asistentes en Maude?\\[8pt]
\rule{\linewidth}{0.4pt}\\[10pt]
\rule{\linewidth}{0.4pt}

\end{enumerate}

\bloque{G. Para terminar}

\begin{enumerate}[leftmargin=*,resume]

\item ¿Qu\'e es lo mejor del proyecto? ¿Qu\'e mantendr\'ias?\\[8pt]
\rule{\linewidth}{0.4pt}\\[10pt]
\rule{\linewidth}{0.4pt}

\item ¿Qu\'e es lo peor? ¿Qu\'e cambiar\'ias?\\[8pt]
\rule{\linewidth}{0.4pt}\\[10pt]
\rule{\linewidth}{0.4pt}

\item ¿Qu\'e consejo le dar\'ias a alguien que empiece este proyecto el curso
que viene?\\[8pt]
\rule{\linewidth}{0.4pt}\\[10pt]
\rule{\linewidth}{0.4pt}

\item \cb\ \textbf{No} quiero que mis respuestas se usen en una publicaci\'on
sobre el diseño de la asignatura.

\end{enumerate}

\bigskip
\begin{center}
\noindent\hrulefill\\
\textit{Gracias por el tiempo dedicado.}
\end{center}

\end{document}
